\section{Case B Reserved deployment and schema migration}\label{sec:caseB}

\subsection{Parameter family and exact owned interfaces}
Fix a finite dependency DAG $(V,\prec)$, finite reservation slots $R$,
a distinct coordinator token $q$, capacity $K\in\mathbb N$ and
nonnegative demands $w:V\to\mathbb N$. Fix $m\geq2$ and the data set
$D=\{0,1\}^m$. Choose a permutation $\pi$ of $\{0,\ldots,m-1\}$ and
mask $b\in D$, and define the bijective schema transformation
\[
h(x)_j=x_{\pi(j)}\mathbin{\mathsf{xor}}b_j.
\]
Its inverse applies the inverse permutation and the corresponding mask.
The two schema versions store physical record $y=x$ at version zero
and $y=h(x)$ at version one.

A registry is an injective partial map $\alpha:V\rightharpoonup R$,
with $L=\dom\alpha$ and version map $s:L\to\{0,1\}$. Legal public
registry tags have predecessor-closed $L$ and
$\sum_{v\in L}w(v)\leq K$. The finite public configuration contains
$\alpha,s$ and any explicitly returned ordinary response. The owned
bundle contains $q:\mathsf{Coordinator}$ and every $r\in R$:
unused slots have kind $\mathsf{Reserved}$ and $r=\alpha(v)$ has kind
$\mathsf{Live}(v,s(v))$. The coordinator has a private registry copy;
a reserved slot has singleton private state; each live slot has a private
payload. Public/private registry agreement and payload validity form the
reachable-state invariant.

Every method takes the \emph{entire} bundle $\{q\}\cup R$ and explicitly
returns all its entries. Unmodified entries are part of the table, not
an implicit linear frame. No identity disappears in this family.
Unrelated resources can be split off only by the ordinary explicit
ownership discipline. This is a serialized finite model, not a scheduler
or concurrency theorem. All registry/version tags and payloads are finite,
so these are genuine O4 dependent interfaces. There is no rollback or
recovery premise.

\subsection{Two genuine physical representations}
Call the record reference implementation $\mathsf{Rec}$ and the packed
candidate $\mathsf{Pack}$. The latter stores
\begin{equation}\label{eq:B-codec}
C(y)=(n(y),p(y)),\quad n(y)=\sum_{j=0}^{m-1}2^jy_j,\quad
p(y)=\sum_jy_j\pmod2.
\end{equation}
Its carrier is the valid subset
\[
P_m=\{(n,p):0\leq n<2^m,\ p=\operatorname{popcount}(n)\pmod2\}.
\]
Define $U(n,p)_j=\lfloor n/2^j\rfloor\pmod2$. Binary expansion gives
$UC=\id_D$ and, on precisely $P_m$, $CU=\id_{P_m}$. An $m$-field record
has become a packed coordinate with a checked redundant bit, not a
renaming of a domain label. Invalid parity states are excluded explicitly.
Logical decoding is $d_0(y)=y$, $d_1(y)=h^{-1}(y)$ for records, and
$d_s(U(z))$ for packed payloads.

Let $F$ map each full record endpoint to its packed endpoint by applying
$C$ at each live slot and leaving public configuration, versions,
coordinator, logical responses and actual token identities/kinds fixed.
Its inverse applies $U$. This is an actual isomorphism between different
expanded endpoint sets. It is not an executable public operation exposing
private store or retyping a capability. At a separately chosen raw-export
presentation, ordinary exported payload endpoints would instead be
transported by $C,U$; those observations are not silently included in the
common hidden-state client interface below.

\subsection{Complete deterministic primitive tables}
Every primitive has one witness of probability one and one output hole.
Public unary events are $\mathsf{Alloc}(v,r)$, $\mathsf{Read}(v)$,
$\mathsf{Write}(v)$, $\mathsf{Migrate}(v)$ and $\mathsf{Reject}(op)$.
Payload is public only when explicitly returned. All operations first check that the private coordinator registry equals
the public registry and that the declared consistency conditions hold.
A failed invariant check emits $\mathsf{RejectInvariant}$ at unit cost,
returns an error flag and preserves the entire store and bundle on both
sides. Otherwise they use the tables below. A failed ordinary enabling
predicate likewise emits its rejection, returns an error flag and preserves
the full store and bundle.
\begin{itemize}
\item $\mathsf{alloc}(v,r,x)$ requires $v\notin L$, unused $r$, all
  predecessors of $v$ in $L$, and
  $\sum_{u\in L}w(u)+w(v)\leq K$. It sets $\alpha(v)=r$, $s(v)=0$,
  retypes that existing slot to $\mathsf{Live}(v,0)$, installs $x$ or
  $C(x)$, updates the coordinator registry and returns success.
\item $\mathsf{read}(v)$ requires $v\in L$. It returns
  $d_{s(v)}(y)$ or $d_{s(v)}(U(z))$ without changing any entry or kind.
\item $\mathsf{write}(v,x)$ requires $v\in L$. It installs
  $h^{s(v)}(x)$ or $C(h^{s(v)}(x))$, leaving registry and kinds unchanged.
\item $\mathsf{migrate}(v)$ requires $v\in L$ and $s(v)=0$.
  It changes $y\mapsto h(y)$ or $z\mapsto C(h(U(z)))$, sets $s(v)=1$,
  updates the coordinator and retypes the slot to $\mathsf{Live}(v,1)$.
  The logical payload is preserved because $h^{-1}h=\id$.
\end{itemize}
All other entries pass through explicitly. The ordinary enabling guards are finite public
predicates, so enabled forms can use O4's singleton/empty evidence, while
the checked forms define arbitrary admitted clients on failed requests.
The success theorem below does not relabel rejection as success. Migration
is a declared effectful state transition with event and positive cost;
purely calculating $h(y)$ cannot retype an owned capability. There is no
release or hidden reuse of slots.

A concrete scalar cost assignment is
\begin{center}
\begin{tabular}{@{}lcc@{}}
\toprule Method&Record reference&Packed candidate\\\midrule
Allocation&3&2\\
Read&2&1\\
Write&3&2\\
Migration&$m+2$&$m+1$\\
Rejection&1&1\\
\bottomrule
\end{tabular}
\end{center}
More generally choose any finite nonnegative vector annotations satisfying
$c_{\mathsf{Pack}}(Fz,M)\leq c_{\mathsf{Rec}}(z,M)$ at conjugate inputs,
including equality. They are supplied model costs, not measured hardware
speed. If cost becomes ordinary client-visible data, equality of that
public datum would additionally be required.

\subsection{Local conjugacy and substantive endpoint squares}
For each method $M$ with its finite public argument, denote its full
state/output map by $M_{\mathsf{Rec}}$ or $M_{\mathsf{Pack}}$.
\begin{proposition}[Schema and representation conjugacy]\label{prop:B-conjugacy}
At every enabled or rejected related input,
\[
F_GM_{\mathsf{Rec}}=M_{\mathsf{Pack}}F_E,
\qquad
M_{\mathsf{Pack}}=F_GM_{\mathsf{Rec}}F_E^{-1},
\]
where the actual source and target endpoint maps apply $C$ on the slots
specified by their registry/version tags. Public responses and capability
kinds are included in the equation.
\end{proposition}
\begin{proof}
Allocation installs $C(x)$ in the packed carrier. Read uses $UC=\id$;
write uses $C$ after the same version encoder. Migration is the exact
identity $Ch=(ChU)C$. The invariant guard agrees because $F$ preserves the private coordinator
copy and public registry. If it fails, both sides preserve their related
full states and emit the same rejection. Otherwise ordinary guards depend
only on public registry, versions, capacity, demands and identities fixed
by $F$, so ordinary rejection also agrees.
Every other entry passes through, and both coordinator updates use the
same new public registry. Valid packed carriers and reachable registry
agreement are preserved in each case.
\end{proof}
The endpoint maps are defined tagwise on finite coproducts, coordinatewise
on owned stores, compose as codecs, and have the displayed inverse.
Ordinary public substitution commutes with them because the fixed private
codec ignores the public substitution environment. They also commute with
products on explicitly disjoint bundles. This supplies substantive
morphism-of-examples data, not merely agreement of example names.

Each method has the actual mixed body
$\theta_i(z)=\mathsf{event}_M(\leaf M_i(z))$.
Conjugacy gives $TF_G\theta_{\mathsf{Rec}}=
\theta_{\mathsf{Pack}}F_E$. The singleton witness map, probability one
and declared cost inequality therefore give a genuine native endpoint
square by~\eqref{eq:native-square}. It is cost-exact for equal annotations
and a noninvertible cost-lowering cell for strict annotations. The
schema transition and the representation endpoint square are distinct:
the former is an operation, the latter is mathematical transport of its
whole interpretation.

\subsection{Full O1 through O5 interpretation}
This case uses precisely the independently verified base, joint rig,
terminal causal-tree content and mixed occurrence construction of
Section~\ref{sec:core}. It selects singleton probability-one witnesses.
That deterministic computation subfragment is closed under the relevant
finite compositions; it is \emph{not} a rig embedding $c\mapsto(1,c)$,
which would fail additive preservation because $1+1\ne1$.
No stochastic assumption or recovery premise is needed for the deployment
or conjugacy theorem.

The O2 action is ordinary-data framing, with
$F_Z=\id_Z\times F$ commuting with the framed body at every leaf.
It does not frame the owned bundle as freely duplicable data. All O4
primitive output sets equal the input identity set at their one leaf,
so the disjoint-subset certificate holds literally. Nested bind, case,
explicit disjoint batch and decreasing-rank invocation provide the
language interpretation. If a topological ordering is $v_1,\ldots,v_k$
and distinct reserve choices are $r_1,\ldots,r_k$, its deployment term
is the finite nested bind of
$\mathsf{alloc}(v_i,r_i,x_i)$, with finite success/error cases and a final
return. Feasibility below proves its success branches are selected when
the stated assumptions hold. Reads, writes and migrations can be appended
as finite typed terms.

Translate method symbols from record to packed while preserving public
pure tables, code tags, captures and constructors. Induction on syntax,
with decreasing-rank induction at invocation, transports the full tree
and lowers accumulated cost: at bind the actual related leaves are
substituted and $TF$ commutes with grafting; at batch use the disjoint
partition; at case equal public data choose the same branch. O3 therefore
gives equal transported joint laws for equal costs and monotone
couplings for optimized costs. Deterministic clients have Dirac laws;
the finite-discrete QBS embedding preserves these comparisons.

For O5, use the full tokenwise relation rather than restricting it by a
cross-token reachable-state condition. The invariant-check completion
above handles all equal private coordinator copies, including copies
that disagree with the public registry. On those inputs both sides reject
identically and preserve their related states. On invariant inputs the
conjugacy calculation applies; payload carriers remain the valid carriers
for their actual kinds. Hence the local certificate applies to every
related full input. Public/private agreement is a separately initialized
and preserved reachable invariant used for deployment safety, not an
unproved extra restriction on O5's relation.

At the common logical interface, use equal coordinator registries,
$*\sim*$ on reserves, and $y\sim C(y)$ at each live kind. All local
method matrices have one entry of mass one on the conjugate outputs,
with equal public node, related private leaves and candidate cost no
larger. The initial empty-registry stores are related. Other fixed client
primitives must preserve this relation, as required by O5; identical
names alone do not establish that premise. Consequently
\begin{equation}\label{eq:B-contextual}
\Obs(C[\mathsf{Pack}])\sqsubseteq\Obs(C[\mathsf{Rec}])
\end{equation}
for every admitted bounded typed client, including adaptive public
choices, repeated calls, permitted ranked callbacks and certified
probabilistic client primitives. This does not quantify over physical
private-store inspection, arbitrary reentry, unbounded recursion or
observable timing. For a physical-export interface instead, the client
must be transported along the actual endpoint codec; different public
observations cannot be identified without that transport.

The causal content is also substantive within this full family. On a
live slot, one read followed by ordinary reply disposal and two reads
followed by reply disposal have the same final store and one final leaf,
but bodies $\mathsf{Read}(v)(\leaf)$ and
$\mathsf{Read}(v)(\mathsf{Read}(v)(\leaf))$. Their List arity shadows
coincide while the public causal histories do not. Native cost reduction
compares matched trees; it does not justify deleting any computation
with the same final payload.

\begin{theorem}[Exact reserved deployment feasibility]\label{thm:B-capacity}
Every successful prefix preserves registry injectivity, predecessor closure,
the initial actual identities and capacity. Starting with empty registry,
a workflow deploying every vertex exists exactly when
\[
|V|\leq|R|\qquad\hbox{and}\qquad\sum_{v\in V}w(v)\leq K.
\]
Enabled finite reads, writes and migrations preserve feasibility and the
schema/representation invariants.
\end{theorem}
\begin{proof}
Induct on primitive steps. Allocation adds an absent vertex and unused
slot only after checking predecessors and new demand. Other methods and
rejections do not change the live set or demand, and the preceding
conjugacy calculation proves their payload invariant. Necessity follows
from the final injection and capacity. For sufficiency choose a
topological ordering and distinct reserve slots. At each step all
predecessors are already live; nonnegative demands make every prefix
sum no larger than the total, so every check succeeds. The empty graph
also works, including with no slots. Migration changes neither demand
nor logical data, and enabled reads/writes preserve the registry.
\end{proof}
For the chain $a\prec b\prec c$ with unit demands, three slots and
$K=2$, the third allocation is rejected. More slots cannot solve that
capacity failure. With two slots and $K=3$, total demand fits but the
third distinct slot is missing. An out-of-order allocation is rejected
even with enough capacity and slots. The theorem concerns the finite
DAG/no-release model, not cyclic dependencies or arbitrary release/reuse
scheduling.

\subsection{Worked schema calculation and a lossy boundary}
Choose $m=3$, $\pi(j)=j+1\pmod3$, $b=(1,0,0)$, edges
$a\prec c$, $b\prec c$, three slots, demands $(1,2,1)$ and capacity
four. Deploy all three vertices with logical $x=(1,0,1)$.
Then
\[
C(x)=(5,0),\qquad h(x)=(1,1,1),\qquad C(h(x))=(7,1).
\]
Migration changes the physical representation but logical read still
returns $(1,0,1)$. Writing $(0,1,0)$ in version one stores
$h(0,1,0)=(0,0,0)$ and packed value $(0,0)$; logical read returns
$(0,1,0)$. Three allocations, one migration and one read cost
$3\cdot3+5+2=16$ for records and $3\cdot2+4+1=11$ for packing.
Nodes, logical results and token kinds agree under the codec; physical
payloads are not literally equal.

A lossy map $L(y)=(y_0,\ldots,y_{m-2})$ cannot replace $C$ at this
interface. Two version-zero states differing only in the last bit have
the same image under $L$ but different permitted read results. A target
read $g$ satisfying $gL=\mathsf{readLast}$ would have to return both
bits from one state. Thus neither this conjugacy nor equal-public-output
O5 certification is available. Restricting the observation language
would be a new boundary, not a proof of the present claim.

\texttt{check\_case\_B\_F.py} checks 124 codec payload states for
$m=2,\ldots,6$, both versions, 10,912 write squares, 320,152 finite
DAG/slot/capacity instances, the worked transitions, the lossy collision and 1,984 off-invariant
method rows for the common $\mathsf{RejectInvariant}$ completion.
These finite checks corroborate the symbolic inverse, tablewise
conjugacy and feasibility proofs; they do not prove a general theorem
by enumeration. All costs, whole-bundle serialization, valid-payload
restrictions, finite/ranked syntax and permitted observations are explicit
model assumptions, not guarantees about an external deployment system.

